\section{Estado del arte}

Históricamente filósofos, psicólogos, sociólogos y pedagogos, han resaltado la importancia de las competencias transversales y socioemocionales (CTSE) en la educación de los individuos, que debe abarcar múltiples dimensiones, incluyendo el desarrollo cognitivo y el socioemocional. Las habilidades “blandas” o “no cognitivas” son rasgos de la personalidad de un individuo que no dependen de conocimientos técnicos previos y que le permite interactuar, comunicarse y trabajar con otros.
\\
\\
La Corporación Andina de Fomento (CAF), el Banco Mundial, el Banco Interamericano de Desarrollo (BID) y la Organización de Cooperación y Desarrollo Económico (OCDE) han sido actores importantes en la promoción y discusiones políticas contemporáneas en torno al fomento de las CTSE. La publicación del reporte “Desconectados” del BID en 2012 aporta una perspectiva muy pertinente a la discusión sobre la calidad de la educación en América Latina, indicando la existencia de un desacople entre las habilidades que se forman actualmente en el sistema educativo y las que son requeridas en el mercado laboral. De igual forma, el reporte de “Más Habilidades para el Trabajo y la Vida” de la CAF en 2016 señala que América Latina necesita incrementar sus esfuerzos en desarrollar en los jóvenes un conjunto de habilidades que les permita navegar y adaptarse a los entornos cambiantes de los próximos años.
\\
\\
La evidencia a nivel internacional basada principalmente en estudios empíricos y longitudinales en países desarrollados destaca la relevancia de las competencias transversales y socioemocionales (CTSE) en el proceso de aprendizaje y en los logros tanto actuales como futuros de los individuos. Por ejemplo, se ha observado una correlación positiva entre estas competencias y la finalización de estudios universitarios, el nivel de ingresos, el bienestar físico y mental, así como con el bienestar subjetivo, entre otros aspectos \cite{Huerta2019}. 
\\
\\
No existe consenso aún frente a cómo nombrar las habilidades blandas y una definición única y universal, dado a que cada entidad o institución tiende a enfocarse en aspectos específicos dentro de este amplio espectro de competencias. Este desafío se evidencia en la diversidad de enfoques y terminología utilizada por diferentes estudios, organizaciones y sectores. Mientras algunas pueden referirse a estas habilidades como competencias transversales y socioemocionales, otras pueden preferir términos como habilidades interpersonales, habilidades para la vida o habilidades no cognitivas \cite{RomeroGonzalez2021}.
\\
\\
Las habilidades blandas contemplan, entre otras, la capacidad de comunicarse, trabajar en equipo, empatía, adaptabilidad, resiliencia, autogestión, proactividad, autocrítica, honestidad, pensamiento analítico e innovación, aprendizaje activo y estrategias de aprendizaje, creatividad, originalidad e iniciativa, pensamiento crítico, liderazgo e influencia social, inteligencia emocional, razonamiento y resolución de problemas. Estas contribuyen al crecimiento personal y profesional del estudiante, así como a su convivencia entre pares y la comunicación efectiva dentro y fuera de una organización \cite{CepedaGuerra2021}. 
\\
\\
En un contexto donde la tecnología avanza y se reinventa rápidamente, el entorno se torna cada vez más competitivo. Adecco \cite{Adecco2024}, la empresa consultora en recursos humanos y de procesos de selección de personal, expone que quiénes desarrollan sus habilidades blandas, son los que obtienen ventaja competitiva frente a los cambios del entorno, ya que la demanda del mercado está apostando por los candidatos que demuestren estas habilidades, incluso las empresas están invirtiendo en programas de desarrollo que fortalezcan estos atributos en sus empleados. 
\\
\\
En el trabajo de grado de “Habilidades blandas y su influencia en el clima laboral de las organizaciones: una mirada actual para Colombia”, se aborda la creciente importancia de las habilidades blandas en el contexto laboral contemporáneo y como influyen de algún modo en el clima laboral de las organizaciones, puesto que contribuyen a la sostenibilidad y productividad de estas. Las empresas buscan ser competitivas no solo con personal técnicamente preparado, sino también con líderes y equipos de trabajo eficaces, eficientes y comprometidos \cite{PerdomoRamos2023}.
\\
\\
Es crucial que las instituciones educativas sean capaces de adaptarse y proporcionar a sus estudiantes alternativas para que adquieran las habilidades necesarias y se conviertan en profesionales altamente cualificados en sus respectivos campos. En la actualidad quiénes responden estas demandas, lo hacen con plataformas tecnológicas de enseñanza y aprendizaje que sirven de complemento al modelo tradicional de aula. 
\\
\\
En el artículo “Praxis sobre el desarrollo de competencias blandas en estudiantes de educación superior en aulas virtuales”, realizan una revisión bibliográfica sobre el desarrollo de competencias blandas en ambientes virtuales de enseñanza superior, comprobando así que las instituciones que se están adaptando a estas necesidades desarrollando plataformas y estrategias, monitorean constantemente el desarrollo de competencias de los estudiantes, evaluando sus avances y estableciendo un plan de mejora continua conforme a objetivos planteados y de la retroalimentación de los estudiantes \cite{CepedaGuerra2021}.
\\
\\
En la investigación “Soft skills in software engineering in higher education. An initial study of its state in Latin America”, se mencionan diversos autores que han discutido sobre la necesidad y relevancia de incorporar en los currículos de los programas de estudio de ingeniería de sistemas la enseñanza o la visibilización sobre habilidades blandas. Plantean que en la actualidad los estudiantes no solo deben desarrollar competencias técnicas, sino también competencias transversales para que realicen actividades profesionales de manera efectiva en un mundo globalizado, sin embargo, son el eslabón perdido en la ingeniería de sistemas, ya que los planes de estudios suelen marginar estas habilidades.
\\
\\
Adicionalmente, en el artículo de investigación anteriormente mencionado, se presentan los resultados de una encuesta realizada a directores, coordinadores y responsables académicos de 115 programas universitarios de ingeniería de sistemas en 85 universidades de 15 países de América Latina sobre el estado actual de la enseñanza de habilidades blandas en estas carreras. Casi la mitad de los programas encuestados incluyen cursos específicos sobre estas habilidades, como comunicación, trabajo en equipo, liderazgo, ética profesional y gestión de conflictos. Sin embargo, hay otras habilidades consideradas relevantes para la ingeniería de software que no se mencionaron, como habilidades de presentación, negociación, resolución de problemas, interés por aprender y autonomía.